\chapter{Conclusion}
\label{conclusion}

\minitoc{}

This manuscript addressed the modelling and control of chains of tethered underwater robots, with a
focus on fast, analyzable cable representations and chain-level controllers that respect tether
geometry. Two technical threads were developed in tandem: (i) cable modelling suitable for
real-time use in control, and (ii) coordination of robot chains under geometry-consistent
constraints. Below, we synthesize the main outcomes, discuss limitations, and outline future work.

\section{Cable Modelling}
\label{c_cable_modelling}

\subsection{\catenarymodeltitle{}}
\label{c_cm_catenary_model}

The classical catenary model was revisited and tooled for control applications
(\Cref{cable_modelling}). In particular:
\begin{itemize}
	\item A compact parameterization was derived that relates the curve to endpoint geometry and cable
	      length. Curvilinear discretization was introduced to generate uniformly spaced samples along the
	      arc, matching how measured points are typically acquired in practice.
	\item The model provides closed-form expressions for horizontal and vertical tensions along the cable,
	      which are directly reusable as external forces in vehicle dynamics.
	\item As expected, the vertical catenary is a strong baseline in quasi-static regimes; its analytical
	      form makes it compatible with real-time control pipelines and optimization-based controllers.
\end{itemize}

\subsection{\augmentedmodeltitle{}}
\label{c_cm_augmented_catenary_model_dynamics}

To accommodate steady hydrodynamic effects under motion, an augmented two-\emph{DOF} extension was
proposed (\Cref{cm_augmented_catenary_model}). The cable curve is obtained by applying two
rotations to the vertical catenary: a sway-related rotation \(\gamma\) around the end-to-end axis
and a surge-related tilt \(\theta\) around the plane normal. This yields the hetagammamodeltitle{}
while preserving analytical tractability and the curvilinear discretization.

The model was evaluated on a motion-capture dataset comprising \num{129} sequences and \num{8}
distinct cables spanning different buoyancies and flexibilities (\Cref{experiments}). The residual
between measured and modelled point clouds, computed frame-wise over uniformly spaced samples, was
used as the accuracy metric. Across cables and conditions, the augmented form systematically
outperformed the purely vertical catenary, with median residual reductions in the
\SIrange{10}{75}{\percent} range; for the most compliant and negatively buoyant cables, typical
median residuals reached \SI{10}{\milli\meter} to \SI{13}{\milli\meter}.

Two avenues toward dynamics for \(\theta\) and \(\gamma\) were explored: a physics-inspired
surrogate (pendulum analogy driven by endpoint motion) and data-driven symbolic regression. While
promising, these preliminary efforts did not yet yield a robust dynamic law suitable for deployment
in chain control. Consequently, the control chapter deliberately used the vertical catenary, for
which simple, reliable tension expressions are available.

\paragraph{Takeaways}
\begin{itemize}
	\item The \catenarymodeltitle{} remains an excellent real-time baseline for quasi-static and mild dynamic
	      regimes.
	\item The \thetagammamodeltitle{} captures steady hydrodynamic deformations with small computational cost
	      and statistically significant accuracy gains across a diverse cable set.
	\item Establishing consistent, well-behaved dynamics for the augmentation parameters is the key remaining
	      step to integrate the \augmentedmodeltitle{} in predictive controllers.
\end{itemize}

\section{Robot Chain Control}
\label{c_robot_chain_control}

Building on a standard marine-craft model and catenary-based tether forces, a unified chain model
was formulated (\Cref{robot_chain_control}). This representation embeds horizontal and vertical
forces from the catenary directly into each vehicle's dynamics, enabling constraints on seafloor
and surface clearance, inter-robot distances, and admissible speeds/thrusts to be stated in a
geometry-consistent fashion.

Three complementary control strategies were examined under common scenarios (varying bathymetry,
leader manoeuvres, mixed surface/underwater vehicles):
\begin{itemize}
	\item A constrained \ac{mpc} that optimizes inputs over a finite horizon while enforcing chain dynamics
	      and catenary-induced constraints.
	\item A planning layer that generates feasible waypoint sequences followed by a \ac{pid} tracker,
	      offering a low-complexity alternative with explicit feasibility filtering.
	\item A lightweight reactive behavior adapted from leader-follower visual servoing, here using relative
	      pose feedback rather than image features.
\end{itemize}

Simulations compared tracking errors, inter-robot spacing, seafloor clearances, and computation
times. The \ac{mpc} offers the most systematic constraint handling and best overall performance
when a sufficient computation budget is available; the planner-tracker strikes a balance between
feasibility and simplicity; reactive control remains attractive for fast safety overrides and as a
fallback layer.

Experiments validated the end-to-end integration with motion-capture feedback (\Cref{experiments}).
We demonstrated: (i) single-robot closed-loop control using a \ac{pd} and then an \ac{mpc} with a
short horizon at \SI{10}{\hertz}; and (ii) a leader-follower setup where the follower solves an
\ac{mpc} using the leader's predicted motion and catenary-derived constraints. Robust operation
required state-estimation fusion via an \ac{ekf} to bridge tracking gaps; practical procedures were
also implemented to recover from solver non-convergence by adaptively relaxing tolerances.

\paragraph{Takeaways}
\begin{itemize}
	\item A compact, unified robot-catenary model supports constraint-aware control design for chains.
	\item Controllers with different computational footprints cover complementary parts of the design space:
	      systematic constrained optimization, feasible planning with simple tracking, and reactive safety
	      behaviors.
	\item Experimental feasibility was shown on Bluerov platforms with motion-capture feedback; additional
	      integration and sensing are needed for field deployment.
\end{itemize}

\section{Limitations}
\label{c_limitations}

The main limitations identified throughout the work are:
\begin{itemize}
	\item \textbf{Augmented parameters dynamics.} A consistent dynamic law for \(\theta\) and \(\gamma\) remains
	      to be established. Without it, the \augmentedmodeltitle{} is best used as a quasi-steady approximation
	      or for frame-wise estimation.
	\item \textbf{Hydrodynamic simplifications.} Drag and added-mass effects were handled implicitly through
	      the augmentation; unsteady flows, strong accelerations, and vortex-induced phenomena are not captured.
	\item \textbf{Chain experiments at scale.} Experimental validation was limited to a pool with one and two
	      Bluerovs and relied on motion capture. Broader trials with longer cables, currents, and real bathymetry
	      are needed.
	\item \textbf{Sensing and estimation.} Robust field operation will require integration of \ac{dvl}, depth,
	      \ac{usbl}/acoustics, and possibly vision/sonar, with reliable multi-rate fusion beyond the \ac{ekf}
	      used here.
	\item \textbf{Computational budget.} Real-time \ac{mpc} in multi-robot configurations remains sensitive to
	      model size, horizon length, and solver settings; careful warm-starting, constraint tightening, and
	      tailored cost design are important in practice.
\end{itemize}

\section{Perspectives and Future Work}
\label{c_perspectives}

The following directions are natural continuations of this thesis:
\begin{itemize}
	\item \textbf{Dynamics and estimation for the \augmentedmodeltitle{}.} Derive and validate dynamic equations
	      for \(\theta\) and \(\gamma\) combining physics priors (e.g., reduced-order pendulum-like models) and
	      learning-based identification from onboard sensing. Study observability and propose filters to estimate
	      the augmented state online.
	\item \textbf{From geometry to mechanics.} Couple the \augmentedmodeltitle{} with lumped-mass or efficient finite
	      elements when needed, switching models depending on regime while keeping interfaces stable for control.
	\item \textbf{Distributed and scalable control.} Investigate distributed \ac{mpc} and hierarchical planning for
	      longer chains, including constraint tightening, invariant sets for safety, and event-triggered reactive
	      layers.
	\item \textbf{Field validation.} Transition from motion capture to onboard sensing (\ac{dvl}, acoustics, vision,
	      pressure) and validate in tanks with current generators and then at sea, including mixed \ac{rov}/\ac{usv}
	      chains and realistic tasks.
	\item \textbf{Safety and interaction with the environment.} Integrate seabed/structure interaction models,
	      explicit clearance guarantees, and contingency handling for slack cables and contact.
\end{itemize}

\section*{Closing Remarks}

This thesis contributes an analytically tractable modelling layer for tethers, an augmented variant
that captures steady hydrodynamic deformations, and a set of chain-control strategies that respect
tether geometry. Together with the experimental demonstrations, these results support the viability
of cable-aware control for robot chains and provide a practical foundation for future
implementation at scale and in the field.
